% ============================================================
\chapter{Introduction}
\label{chap:introduction}
% ============================================================


% -------------------------------------------------------
\section{Research Motivation  }
% -------------------------------------------------------

Breast cancer is the most commonly diagnosed malignancy worldwide
and remains a leading cause of cancer-related mortality among women.
Despite substantial advances in early detection and targeted therapy,
the molecular mechanisms governing the earliest phases of tumour
initiation remain incompletely understood.

Carcinogenesis is not a binary transition from normal to malignant,
but a progressive deviation from epigenetic homeostasis.
The epigenetic hallmarks of frank tumour tissue, including
genome-wide hypomethylation and focal hypermethylation of
tumour-suppressor loci, have been extensively characterised
\cite{ehrlich2002,bird2002}. The methylation signatures of
invasive breast carcinoma are therefore well established.
While the existence of epigenetic field defects in Normal-Adjacent
breast tissue has been established \cite{Teschendorff2016},
their reproducible and transferable characterisation across
independent cohorts — under strict statistical control —
remains an open methodological challenge.

This observation forms the basis of the \textit{field cancerisation}
hypothesis, originally proposed in oral cancer and subsequently
extended to breast tissue through genome-wide methylation profiling
\cite{Teschendorff2016}. The field effect posits that carcinogenesis
does not arise from a single isolated cell but emerges within a
broader epigenetically conditioned tissue environment —
a molecularly altered field that remains invisible to the pathologist
yet detectable at the epigenomic level.

DNA methylation is particularly well suited for investigating
such early alterations. Methylation marks are mitotically heritable,
chemically stable, and measurable at single-locus resolution
across hundreds of thousands of CpG sites using Illumina array
platforms. Prior work has demonstrated that alterations
accumulating in Normal-Adjacent breast tissue preferentially
target Polycomb-regulated and bivalent chromatin loci,
precisely the genomic regions implicated in cancer-associated
silencing \cite{Teschendorff2016,Teschendorff2010}.
Critically, these alterations manifest predominantly as
increased methylation \emph{variability} rather than as large
directional mean shifts, rendering classical differential
methylation analyses insufficient for their detection
\cite{Teschendorff2014}.

This thesis addresses whether DNA methylation alone can
reproducibly discriminate Normal from Normal-Adjacent breast
tissue across independent cohorts under strict statistical control.

\section{Unresolved Limitations and Open Problem}

Despite the biological importance of the Normal--Adjacent state,
four unresolved limitations motivate the present work.

\paragraph{Mean-based differential methylation is insufficient.}

Standard approaches identify discriminative CpGs via
comparison of group means using linear models or t-tests.
However, field-effect alterations are expected to manifest
as dispersion increases rather than uniform shifts.
Teschendorff and colleagues formalised the distinction between
differentially methylated (DM) and differentially variable (DV)
loci, demonstrating that DV CpGs in Normal-Adjacent tissue are
enriched for biologically meaningful genomic features
\cite{Teschendorff2014,Teschendorff2016}.
A reliance on mean-based selection therefore risks missing
a substantial component of the pre-malignant signal.

\paragraph{Field defects are heterogeneous and cohort-dependent.}

Epigenetic alterations associated with field cancerisation
are not homogeneous across individuals or genomic loci.
Moreover, as documented in the analyses developed
in Chapters~\ref{chap:feature_selection}
and~\ref{chap:model_training_evaluation}, the direction
of epigenetic drift along the Normal--Adjacent axis
is not consistent across independent cohorts.
Such polarity heterogeneity has direct consequences for
predictive transfer and has not been systematically
characterised in prior work.

\paragraph{High-dimensional methylation data are structurally unstable.}

Illumina 450K and EPIC arrays profile up to $8.6 \times 10^5$
CpG loci simultaneously, while available breast tissue
cohorts typically include fewer than $n=300$ samples.
This high-dimensional, low-sample-size (HDLSS) regime
($p \gg n$) induces instability in marginal feature selection
procedures \cite{Meinshausen2010}.
Small perturbations in the training set may yield markedly
different CpG rankings, particularly in the presence of
strong local correlation among neighbouring probes.
Information leakage during preprocessing or feature
selection can further inflate performance estimates
\cite{Ambroise2002}.

\paragraph{Cross-dataset reproducibility is rarely enforced
under polarity-aware constraints.}

When methylation-based classifiers targeting the
Normal--Adjacent distinction are evaluated across
independent cohorts, the implicit assumption is that
the direction of epigenetic drift is consistent across
studies.
This assumption is rarely tested explicitly.
The present work demonstrates that drift orientation
is not universally preserved: one of the three cohorts
examined exhibits a systematically inverted drift axis
relative to the other two, producing directed
misclassification rather than random degradation
in transfer settings.
Failing to account for this structural heterogeneity
leads to misleading conclusions about signature
transferability.\\



\noindent These four limitations define the central open problem
addressed in this thesis.

\begin{quote}
\textit{Is it possible to construct a robust, transferable
CpG-based signature that discriminates Normal from
Normal-Adjacent breast tissue, while controlling for
high-dimensional instability and enforcing strict
cross-cohort reproducibility?}
\end{quote}


% -------------------------------------------------------
\section{Central Hypothesis}
% -------------------------------------------------------  
 
The central hypothesis of this thesis is that
Normal-Adjacent breast tissue occupies a
quantitatively characterisable intermediate
position along an epigenetic trajectory connecting
histologically normal tissue to invasive tumour.
This Normal--Adjacent--Tumour axis is not merely
conceptual but measurable from high-dimensional
methylation data under appropriate statistical constraints.
This hypothesis yields four testable implications.
First, intra-dataset discrimination of Normal from
Normal-Adjacent tissue should be statistically feasible
using a variance-sensitive, leakage-controlled CpG panel,
yielding AUC substantially above chance.
Second, cohorts exhibiting concordant drift orientation
should support bidirectional cross-dataset transfer
with non-trivial predictive performance.
Third, polarity-inverted cohorts should produce
systematic misclassification in transfer settings,
leading to directed sub-random AUC values rather
than random degradation.
Fourth, a compact CpG panel derived within a single
cohort should retain measurable discriminative power
when transferred to an external cohort with aligned
drift orientation. Failure of these predictions would weaken or refute
the proposed model.

% -------------------------------------------------------
\section{Objectives and Thesis Structure}
% -------------------------------------------------------

The thesis pursues five methodological objectives,
addressed across the following chapters.

Chapter~\ref{chap:biological_background} establishes
the biological and epigenetic background.

Chapter~\ref{sec:data_preparation_exploration} introduces
the three cohorts (GSE69914, GSE225845, GSE287331)
and characterises their intra- and inter-dataset structure.

Chapter~\ref{chap:data_preprocessing} constructs a
leakage-controlled preprocessing pipeline enforcing
strict separation between training and test data
at all stages.

Chapter~\ref{chap:feature_selection} develops a
stability-aware feature selection framework integrating
variance sensitivity, directional stability, correlation
pruning, and biological priors, and reports
inter-dataset concordance analyses.

Chapter~\ref{chap:model_training_evaluation} evaluates
predictive performance under intra- and cross-dataset
designs, derives a compact 50-CpG panel, and
characterises the consequences of drift polarity
heterogeneity for signature transferability.

Finally, Chapter~\ref{chap:multi_cohort_integration} integrates
findings across cohorts and discusses implications
for epigenetic biomarker discovery.